Spectacle de danse : Yé ! (L’Eau), par Circus Baobab, 23 novembre 2024 à la Maison de la culture d’Arlon. 
Dr. P. Jacquemart

\section{Résumé et justification de l'activité}

\todoH{FloraStyle}{red!30}{\ les corps du spectacle  sont différents par rapport à l’europe (choix pe des gens plus filiformes dans les compagnies europiénnes) pe ajouter une ligne sur ce thème}{}

L’activité culturelle que je souhaite détailler est un spectacle de \todoH{FloraStyle}{red!30}{\ est-ce que c'est la bonne façon de citer ou i faut mettre en textcite ?? }{danse, Yé ! (L’Eau)} \parencite{circusbaobab_site}. La danse est un moyen de communication unique soumis l’interprétation de chacun. Cet art requiert une certaine sensibilité et imagination afin d’en comprendre toute son histoire. Apprendre à voir au-delà, comprendre un message au travers des lignes du corps au lieu de celles des livres. La danse est un moyen d’ouvrir l’esprit afin de mieux comprendre l’autre. Tout le monde n’a pas le même langage et pourtant en suivant la gestuelle de la danse, une union se fait entre les spectateurs et les danseurs.  Ayant moi-même pratiqué la danse, je ne pouvais qu’admirer tous ces corps se mouvant à l'unisson. Poids plume, les danseurs étaient impressionnants et au-delà de la prouesse physique, ils ont, selon moi, su transmettre un message fort.
 
Yé ! représente avec force la question de nos ressources actuelles, et leur absence cruelle dans certaines régions du monde. Les danseurs se battent pour atteindre cette eau inaccessible. Ils représentent les Hommes qui se réinventent, construisent mais se déconstruisent également entre eux au nom de la santé, de la vie. Ces voltigeurs créent de véritables pyramides humaines aussi fortes que fragiles grâce à un équilibre étroit. A travers ce spectacle, les danseurs tentent de nous transporter et de nous transmettre l’espoir d’un avenir meilleur. Si, ensemble, nous agissons malgré une situation actuelle plus qu’alarmante. La thématique environnementale était rythmée par la friction des bouteilles en plastique, symboles percutants de la crise climatique actuelle et le fardeau de nos déchets. En effet, le décor était plus que minimaliste, un tapis de bouteilles était l’unique accessoire des danseurs. Loin d’un ballet classique, le contemporain permet une plus grande liberté d’expression et de mouvements. Une empreinte circassienne a également imprégné cette performance.
 
L’un des points de vue inspirants de ce spectacle est la place qu’il accorde à ses danseurs en tant qu’individus et artistes. En effet, le Circus Baobab est une troupe de cirque itinérante. Il a vu le jour en Guinée grâce à Laurent Chevallier, un célèbre réalisateur de cinéma. La compagnie permet à de nombreux danseurs d’avoir un métier stable et de sortir de la précarité. Alors que leur tournée ne s’étendait qu’aux villages environnants, de nos jours, ils traversent toute la France et ses pays limitrophes pour s’exprimer et faire voyager les spectateurs à travers la danse. Ils utilisent les réseaux sociaux pour faire connaître leur cause. Ils n’hésitent pas à montrer leur engagement auprès des plus jeunes car ils recrutent régulièrement. Ils publient également les heures d’entraînement passées afin de nous faire vivre une expérience hors norme ainsi que leur succès et toutes les différentes populations qu’ils ont réussi à toucher avec leurs messages. En outre, ces représentations révèlent des silhouettes rarement observées au sein des compagnies européennes, elles constituent ainsi une invitation à la tolérance et à l'acceptation de la pluralité des corps.

Lorsque le rideau tombe, une réflexion naît. Une fois de plus c’est un rappel de la chance que nous avons dans notre pays, les ressources et le système de santé aussi développés tout en ajoutant une touche de culpabilité car si d’autres régions sont moins chanceuses, des actions peuvent être faites chez nous pour changer les choses. Avoir un rappel aussi délicat et artistique permet d’écouter le message sans se sentir impuissant. De plus, sachant l’histoire de la compagnie, se dire qu’en regardant simplement un spectacle nous rendons possibles à ces danseurs de continuer leur route et de soutenir leurs actions. Ce spectacle ne demande qu’à être regardé, écouté et à ouvrir la porte à la réflexion.

En dehors du cadre académique, je suis convaincue qu’il est primordial d’être ouvert à d’autres façons d’apprendre et de voir le monde. Spécialement dans notre médecine actuelle qui fait face à bien plus de \todoH{FloraStyle}{red!30}{\ défis ? lesquels frustrations de pas les détailler, changer la formulation pour fermer les propositions  }{défis} que les soins du patient. Parmi ces défis, je peux notamment citer l'élaboration particulière de la relation patient-médecin ou encore la conscience de la santé planétaire ou même les aspects économiques des soins.